% Build on the CVPR author kit template, which supplies cvpr.sty and the bst.
% Numbers come from tables/, generated by experiments/make_tables.py. Never
% type a number into the prose; add a macro.

\documentclass[10pt,twocolumn,letterpaper]{article}

% Must precede cvpr.sty, which reads them at load time.
\def\paperID{0000}
\def\confName{CVPR}
\def\confYear{2026}

% Fallback layout is for drafting only and is not submission accurate.
\IfFileExists{cvpr.sty}{%
  \usepackage{cvpr}%
}{%
  \usepackage[letterpaper,margin=0.8in]{geometry}
  \usepackage{times}
  \setlength{\columnsep}{0.31in}
  \typeout{^^J*** cvpr.sty NOT FOUND: draft layout, not submission format ***^^J}
}

\usepackage{graphicx}
\usepackage{amsmath}
\usepackage{amssymb}
\usepackage{booktabs}
\usepackage{xcolor}

% hyperref must load last.
\definecolor{cvprblue}{rgb}{0.21,0.49,0.74}
\usepackage[pagebackref,breaklinks,colorlinks,allcolors=cvprblue]{hyperref}

% Overleaf uploads land flat; local layout keeps them in tables/.
\IfFileExists{tables/numbers.tex}{% Generated by experiments/make_tables.py. Do not edit by hand.
\newcommand{\bumpHapke}{7.4}
\newcommand{\bumpTilt}{13}
\newcommand{\collapsePredicted}{34.30}
\newcommand{\condCorr}{0.992}
\newcommand{\condErrMax}{74.3}
\newcommand{\condMax}{1335}
\newcommand{\degeneratePixels}{73}
\newcommand{\diligentAgreeCount}{9}
\newcommand{\diligentAgreeDelta}{0.005}
\newcommand{\diligentAvg}{15.37}
\newcommand{\diligentBall}{4.10}
\newcommand{\diligentMaxDelta}{0.156}
\newcommand{\diligentObjects}{10}
\newcommand{\diligentWorstObject}{pot2}
\newcommand{\lightExponent}{-0.50}
\newcommand{\maeCollapse}{34.30}
\newcommand{\maeHapke}{12.98}
\newcommand{\maeLambertian}{7.2e-15}
\newcommand{\noiseExponent}{1.00}
\newcommand{\officialAgree}{0.0003}
\newcommand{\rigCount}{12}
\newcommand{\rigMax}{14.6}
\newcommand{\rigMin}{10.4}
\newcommand{\sphereTilt}{36}
\newcommand{\sweepCond}{2.02}
\newcommand{\sweepLights}{10}
}{% Generated by experiments/make_tables.py. Do not edit by hand.
\newcommand{\bumpHapke}{7.4}
\newcommand{\bumpTilt}{13}
\newcommand{\collapsePredicted}{34.30}
\newcommand{\condCorr}{0.992}
\newcommand{\condErrMax}{74.3}
\newcommand{\condMax}{1335}
\newcommand{\degeneratePixels}{73}
\newcommand{\diligentAgreeCount}{9}
\newcommand{\diligentAgreeDelta}{0.005}
\newcommand{\diligentAvg}{15.37}
\newcommand{\diligentBall}{4.10}
\newcommand{\diligentMaxDelta}{0.156}
\newcommand{\diligentObjects}{10}
\newcommand{\diligentWorstObject}{pot2}
\newcommand{\lightExponent}{-0.50}
\newcommand{\maeCollapse}{34.30}
\newcommand{\maeHapke}{12.98}
\newcommand{\maeLambertian}{7.2e-15}
\newcommand{\noiseExponent}{1.00}
\newcommand{\officialAgree}{0.0003}
\newcommand{\rigCount}{12}
\newcommand{\rigMax}{14.6}
\newcommand{\rigMin}{10.4}
\newcommand{\sphereTilt}{36}
\newcommand{\sweepCond}{2.02}
\newcommand{\sweepLights}{10}
}

\begin{document}

\title{Quantifying the Breakdown of Lambertian Photometric Stereo\\under Minnaert Reflectance}

% TODO: affiliation line before submission.
\author{Nhi Ngo}

\maketitle

\begin{abstract}
Photometric stereo recovers a per pixel surface normal from several images of a
static scene under known, varying illumination, and the standard solver assumes
the surface is Lambertian. Real surfaces are not. We quantify what that
assumption costs. Writing the Minnaert reflectance model in its normalized form
and viewing from a fixed viewpoint, the image stack reduces to a per pixel power
law applied to Lambertian shading. The per pixel factor cancels in the recovered
direction, so the estimated normal is unaffected by albedo, by source
irradiance, and by the surface orientation relative to the camera, while the
estimated albedo carries the full bias. We derive a first order expression for
the angular error, confirm it against a controlled synthetic sweep with two
analytically fixed endpoints, and evaluate the same solver on a benchmark of
real objects with ground truth normals, reproducing the reference baseline to
within \diligentAgreeDelta{} degrees on \diligentAgreeCount{} of
\diligentObjects{} objects. At the simplified Hapke point the error reaches
\rigMin{} to \rigMax{} degrees across the light configurations we test, a range
that overlaps the mid tier objects of that benchmark.
\end{abstract}

\section{Introduction: shape from varying illumination}

Recovering the three dimensional shape of a surface from images is a central
problem in computer vision, and one of the oldest cues for it is shading. When a
surface is lit from a known direction, its apparent brightness varies with the
angle between the local surface orientation and the incident illumination, so
brightness carries information about shape. Photometric stereo makes this
practical by capturing several images of a static scene under different known
light directions and solving for a surface normal at every pixel
\cite{woodham1980}.

The technique matters wherever fine surface geometry has to be recovered without
contact. Industrial inspection uses it to detect surface defects too shallow for
range sensing, and image relighting uses the recovered normals to render a
captured object under illumination it was never photographed in. Both
applications share an assumption that is rarely examined: that the surface
reflects light diffusely, in the Lambertian sense, so that observed brightness
depends on the incident direction but not on the direction from which the
surface is viewed.

Real materials violate this. The question this paper asks is not whether the
assumption fails, which is well known, but by how much and as a function of
what. We answer it by rendering surfaces whose reflectance is controlled by a
single continuous parameter, running an unmodified Lambertian solver on all of
them, and measuring the resulting angular error against ground truth that is
exact by construction.

Our contributions are the following. We show that from a fixed viewpoint the
Minnaert model collapses to a per pixel power law on Lambertian shading, and
that the resulting per pixel factor cancels in the recovered direction while
surviving in full in the recovered albedo. We derive a first order expression
for the angular error and state the range over which it holds. We validate the
solver on the DiLiGenT benchmark, matching the reference normal maps
distributed with it to under \officialAgree{} degrees per object, and we identify a set
of degenerate ground truth pixels in that benchmark that changes one reported
figure by \diligentMaxDelta{} degrees depending on how they are scored.

\section{Related work: non Lambertian photometric stereo}

% TODO: Woodham, shape from shading, the benchmark, the shipped non Lambertian
% methods with their numbers, and learned methods at 5 to 7 degrees. Say that
% inverting a power law is likely known; the claim is diagnosis, not inversion.

\section{Method: image formation and the linear solver}

\subsection{Lambertian image formation}

Let $\hat{n}(x)$ be the unit surface normal at pixel $x$, $\rho(x)$ the albedo,
and $\hat{s}_j$ the unit direction toward the $j$th of $m$ known distant light
sources. Under the Lambertian assumption, and away from shadow, the observed
radiance is
%
\begin{equation}
  I_j(x) \;=\; \rho(x)\,\bigl(\hat{n}(x)\cdot\hat{s}_j\bigr).
  \label{eq:lambert}
\end{equation}
%
This is linear in the product $g = \rho\,\hat{n}$. Stacking the light directions
as the rows of a matrix $L \in \mathbb{R}^{m\times3}$ gives $I = Lg$ at every
pixel, and for $m \geq 3$ light directions that are not coplanar in azimuth the
system is overdetermined and well conditioned.

\subsection{Recovering the normal}

Minimizing the squared residual $\lVert I - Lg \rVert^2$ is an ordinary least
squares problem whose normal equations give
%
\begin{equation}
  g \;=\; (L^{\top}L)^{-1}L^{\top}I,
  \qquad
  \rho = \lVert g \rVert,
  \qquad
  \hat{n} = \frac{g}{\lVert g \rVert}.
  \label{eq:woodham}
\end{equation}
%
The albedo is the length of the recovered vector and the normal is its
direction. The solve is independent at each pixel and is applied to the whole
image as a single matrix product.

\subsection{Minnaert reflectance}

To control how far the surface departs from Lambertian we use the Minnaert
model \cite{minnaert1941} in its normalized form. With $\theta_i$ the incident
angle and $\theta_r$ the viewing angle,
%
\begin{equation}
  f_r \;=\; \frac{c+1}{2\pi}\,k(\lambda)\,
            \bigl(\cos\theta_i\cos\theta_r\bigr)^{c-1},
  \qquad 0 \leq c \leq 1,
  \label{eq:minnaert}
\end{equation}
%
so that the outgoing radiance under a distant source of irradiance $E_0$ is
%
\begin{equation}
  L_r \;=\; \frac{c+1}{2\pi}\,k(\lambda)\,E_0\,
            \cos^{c}\theta_i\,\cos^{c-1}\theta_r .
  \label{eq:minnaert-radiance}
\end{equation}
%
Setting $c = 1$ recovers the Lambertian case, and $c = 1/2$ gives the simplified
Hapke model. We use the simplified Hapke form throughout. The full Hapke model,
which additionally accounts for multiple scattering, shadowing and porosity, is
not a special case of Minnaert, and no claim here extends to it.

\subsection{Reduction to a power law}

% TODO: the reduction, P1 and P2, and the first order error law.

\section{Experiments: real objects with ground truth normals}

\subsection{Benchmark and protocol}

We evaluate on the DiLiGenT photometric stereo benchmark, \diligentObjects{}
real objects spanning a wide range of reflectance, from a near diffuse ball to
strongly specular metallic pieces \cite{shi2016}. Each object provides 96
images at $612\times512$ under calibrated directional lighting, together with a
foreground mask and a ground truth normal map obtained independently of the
photometric measurements. The capture is orthographic and single view, which
matches the assumption used in Section 3.

Each observation is divided by the measured intensity of its light source before
the colour channels are combined, following the protocol distributed with the
benchmark. The solver of Equation \eqref{eq:woodham} is applied without
modification, that is, with no changes made for real input relative to the
synthetic experiments.

\subsection{Evaluation}

We report the mean angular error in degrees between the estimated and ground
truth normals over the foreground mask,
%
\begin{equation}
  \varepsilon(x) \;=\;
  \operatorname{atan2}\!\left(
    \lVert \hat{n}_{\text{est}} \times \hat{n}_{\text{gt}} \rVert,\;
    \hat{n}_{\text{est}} \cdot \hat{n}_{\text{gt}}
  \right).
  \label{eq:angular-error}
\end{equation}
%
The two argument form is used rather than the arc cosine of the dot product
because the derivative of $\arccos$ is unbounded at zero angle, which limits its
usable resolution to roughly $10^{-6}$ degrees. That matters here because the
synthetic Lambertian case is exact to machine precision and would otherwise be
reported at the resolution limit of the metric rather than of the solver.

\begin{table}
  \centering
  \IfFileExists{tables/table1_diligent.tex}%
    {% Generated by experiments/make_tables.py. Do not edit by hand.
\begin{tabular}{lrrr}
\toprule
Object & Ours & Reference & $\Delta$ \\
\midrule
ball & 4.10 & 4.10 & -0.004 \\
bear & 8.39 & 8.39 & -0.001 \\
buddha & 14.92 & 14.92 & +0.001 \\
cat & 8.41 & 8.41 & +0.003 \\
cow & 25.60 & 25.60 & -0.001 \\
goblet & 18.50 & 18.50 & -0.000 \\
harvest & 30.62 & 30.62 & +0.005 \\
pot1 & 8.89 & 8.89 & +0.004 \\
pot2 & 14.46 & 14.65 & -0.186 \\
reading & 19.80 & 19.80 & +0.003 \\
\midrule
Average & 15.37 & 15.39 & -0.020 \\
\bottomrule
\end{tabular}
}{% Generated by experiments/make_tables.py. Do not edit by hand.
\begin{tabular}{lrrr}
\toprule
Object & Ours & Reference & $\Delta$ \\
\midrule
ball & 4.10 & 4.10 & -0.004 \\
bear & 8.39 & 8.39 & -0.001 \\
buddha & 14.92 & 14.92 & +0.001 \\
cat & 8.41 & 8.41 & +0.003 \\
cow & 25.60 & 25.60 & -0.001 \\
goblet & 18.50 & 18.50 & -0.000 \\
harvest & 30.62 & 30.62 & +0.005 \\
pot1 & 8.89 & 8.89 & +0.004 \\
pot2 & 14.46 & 14.65 & -0.186 \\
reading & 19.80 & 19.80 & +0.003 \\
\midrule
Average & 15.37 & 15.39 & -0.020 \\
\bottomrule
\end{tabular}
}
  \caption{Mean angular error in degrees on the benchmark, against the
  reference Lambertian baseline distributed with it. \diligentAgreeCount{} of
  \diligentObjects{} objects agree to within \diligentAgreeDelta{} degrees. The
  exception, \diligentWorstObject{}, is discussed in
  Section~\ref{sec:baseline}.}
  \label{tab:diligent}
\end{table}

\subsection{Baseline results}
\label{sec:baseline}

Table \ref{tab:diligent} reports the result on all \diligentObjects{} objects.
The average of \diligentAvg{} degrees agrees with the reference baseline, and
\diligentAgreeCount{} of the \diligentObjects{} objects match to within
\diligentAgreeDelta{} degrees. Beyond the aggregate, the estimated normal maps
agree with the reference normal maps distributed with the benchmark to under
\officialAgree{} degrees per object, which fixes the solver and the data
handling conventions rather than only their summary statistic.

Two details of the protocol are worth recording because they change the reported
figures. The reference pipeline converts colour to grey with a routine that
clamps its output to the unit interval. Without that clamp, saturated pixels
inside specular highlights disagree by up to 16 degrees and the ball error reads
$4.17$ rather than \diligentBall{}. The clamp affects only 0.02 percent of
observations, but those observations concentrate in highlights and lower the
reported error, so it acts as a mild outlier suppressor inside what is presented
as an unmodified Lambertian baseline.

The second detail concerns \diligentWorstObject{}, the one object in
Table~\ref{tab:diligent} that does not agree to within \diligentAgreeDelta{}
degrees. Its ground truth normal map contains 73 pixels of exactly zero length
inside the mask, so the true orientation there is undefined. The reference
metric evaluates $\arccos$ of the dot product and therefore scores those pixels
at 90 degrees, while a magnitude invariant formulation would score them at
zero, that is, as perfect matches. The 73 pixels account for the entire
difference, $73 \times 90 / 35278 = \diligentMaxDelta{}$ degrees. We exclude
them and report the count, rather than allowing undefined pixels to enter the
mean under either convention.

% TODO: synthetic sweep, ablations, effective exponent.

\section{Discussion and limitations}

% TODO: orthographic only; attached shadows only; simplified Hapke only; the
% validity range of the first order law; the oracle exponent as a bound, not a
% method.

\section{Conclusion}

% TODO

{
    \small
    % Fall back to a built in style when the kit's bst is absent.
    \IfFileExists{ieeenat_fullname.bst}{\bibliographystyle{ieeenat_fullname}}%
                                       {\bibliographystyle{ieeetr}}
    \bibliography{main}
}

\end{document}
